\documentclass[11pt,oneside]{article}

\usepackage[T1]{fontenc}
\usepackage{lmodern}
\usepackage{microtype}
\usepackage[margin=1in,headheight=14pt]{geometry}
\usepackage{amsmath,amssymb,amsthm,mathtools}
\usepackage{booktabs,longtable,tabularx,array}
\usepackage{enumitem}
\usepackage{xcolor}
\usepackage{xurl}
\usepackage{hyperref}
\usepackage{fancyhdr}
\usepackage{titlesec}
\usepackage{listings}
\usepackage{needspace}

\definecolor{deepblue}{RGB}{24,75,125}
\definecolor{deepgreen}{RGB}{16,105,65}
\definecolor{amber}{RGB}{151,91,0}
\definecolor{deepred}{RGB}{145,35,42}
\definecolor{softgray}{RGB}{246,247,249}
\definecolor{rulegray}{RGB}{93,101,111}

\hypersetup{
  colorlinks=true,
  linkcolor=deepblue,
  citecolor=deepgreen,
  urlcolor=deepblue,
  pdftitle={The Alaoglu--Erdos Conjecture: Rotation Discriminants, Matched-Gap Energy, and Radical Norm Compression},
  pdfauthor={Independent continuation report},
  pdfsubject={New exact identities and research targets for the two-base integer-exponent conjecture}
}
\urlstyle{same}

\pagestyle{fancy}
\fancyhf{}
\lhead{Alaoglu--Erd\H{o}s continuation}
\rhead{Rotation discriminants and matched gaps}
\cfoot{\thepage}
\renewcommand{\headrulewidth}{0.4pt}

\titleformat{\section}{\Large\bfseries}{\thesection}{0.75em}{}
\titleformat{\subsection}{\large\bfseries}{\thesubsection}{0.75em}{}
\titleformat{\subsubsection}{\normalsize\bfseries}{\thesubsubsection}{0.75em}{}

\setlength{\parindent}{0pt}
\setlength{\parskip}{0.55em}
\setlength{\emergencystretch}{3em}
\setlist[itemize]{leftmargin=2em,itemsep=0.28em,topsep=0.4em}
\setlist[enumerate]{leftmargin=2.25em,itemsep=0.28em,topsep=0.4em}
\allowdisplaybreaks

\newtheorem{theorem}{Theorem}[section]
\newtheorem{proposition}[theorem]{Proposition}
\newtheorem{lemma}[theorem]{Lemma}
\newtheorem{corollary}[theorem]{Corollary}
\newtheorem{conjecture}[theorem]{Conjecture}
\newtheorem{criterion}[theorem]{Criterion}
\theoremstyle{definition}
\newtheorem{definition}[theorem]{Definition}
\newtheorem{question}[theorem]{Research question}
\theoremstyle{remark}
\newtheorem{remark}[theorem]{Remark}
\newtheorem{warning}[theorem]{Audit warning}

\newcommand{\N}{\mathbb N}
\newcommand{\Nzero}{\mathbb N_0}
\newcommand{\Z}{\mathbb Z}
\newcommand{\Q}{\mathbb Q}
\newcommand{\R}{\mathbb R}
\newcommand{\C}{\mathbb C}
\newcommand{\floor}[1]{\left\lfloor #1\right\rfloor}
\newcommand{\ceil}[1]{\left\lceil #1\right\rceil}
\newcommand{\fract}[1]{\left\{#1\right\}}
\newcommand{\norm}[1]{\left\lVert #1\right\rVert}
\newcommand{\Li}{\operatorname{Li}}
\newcommand{\ord}{\operatorname{ord}}
\newcommand{\rad}{\operatorname{rad}}
\newcommand{\vpp}[2]{v_{#1}\!\left(#2\right)}
\newcommand{\thetaq}{\vartheta}
\newcommand{\binomNtwo}{\binom{N}{2}}
\newcommand{\binomNthree}{\binom{N}{3}}

\newcommand{\statusprior}{\textcolor{deepblue}{\textbf{[imported]}}}
\newcommand{\statusnew}{\textcolor{deepgreen}{\textbf{[new proof]}}}
\newcommand{\statuscomp}{\textcolor{amber}{\textbf{[computation]}}}
\newcommand{\statusopen}{\textcolor{deepred}{\textbf{[open target]}}}
\newcommand{\statusaudit}{\textcolor{amber}{\textbf{[public-status audit]}}}

\newenvironment{resultbox}
  {\begin{center}\begin{minipage}{0.93\textwidth}\hrule\vspace{0.65em}}
  {\vspace{0.45em}\hrule\end{minipage}\end{center}}

\lstdefinestyle{pythonstyle}{
  language=Python,
  basicstyle=\ttfamily\footnotesize,
  backgroundcolor=\color{softgray},
  frame=single,
  rulecolor=\color{rulegray},
  breaklines=true,
  showstringspaces=false,
  columns=fullflexible,
  keepspaces=true,
  keywordstyle=\bfseries\color{deepblue},
  commentstyle=\itshape\color{deepgreen},
  numbers=left,
  numberstyle=\tiny\color{rulegray},
  numbersep=7pt
}

\begin{document}

\hypersetup{pageanchor=false}
\begin{titlepage}
\centering
\vspace*{1.3cm}
{\Huge\bfseries The Alaoglu--Erd\H{o}s Conjecture\par}
\vspace{0.35cm}
{\LARGE Rotation Discriminants, Matched-Gap Energy,\par}
{\LARGE and Radical Norm Compression\par}
\vspace{0.7cm}
{\large An independent continuation of the supplied unified report\par}
\vspace{1.0cm}
\begin{minipage}{0.88\textwidth}
\small
\textbf{Scope.}  This document does not repeat the supplied report's twenty-seven
continuations, its Lean inventory, or its proofs of the primitive-generator theorem.  It
imports only the exact conditional free-monoid structure and the finite-type consequence of
Baker's theorem, then develops a different object: the Vandermonde product of the compact
rotation orbit.  The resulting formulas expose an exact second-order imbalance between the
base-$2$ and base-$3$ integer gap products and a concrete cross-prime transport problem.

\textbf{Trust convention.}  Every imported statement is marked \statusprior.  Every theorem
proved here is marked \statusnew.  Numerical values are marked \statuscomp.  No Lean
verification of the new material is claimed; a formalization blueprint is supplied.  The
full Alaoglu--Erd\H{o}s conjecture is not claimed proved.
\end{minipage}
\vfill
{\large August 22, 2026\par}
\end{titlepage}
\hypersetup{pageanchor=true}

\pagenumbering{roman}

\begin{abstract}
Let $\thetaq=\log_2 3$.  The Alaoglu--Erd\H{o}s conjecture says that
$2^x,3^x\in\Z$ implies $x\in\Z$.  Conditional on failure, the supplied report proves that
there is a least nonintegral solution $\beta$ and integers
$M=2^\beta$, $A=3^\beta$ such that every solution is $n+k\beta$.  Writing
$r_k=\floor{k\beta}$ and $t_k=\fract{k\beta}$, the compact orbit is
$(2^{t_k},3^{t_k})=(M^k/2^{r_k},A^k/3^{r_k})$.

This continuation studies the products of all pairwise differences of that orbit.  It proves
an exact lag-by-lag factorization into two Sudler-type factors, an exact structural-prime
denominator formula, an integer orbit-polynomial discriminant identity, and a logarithmic
energy limit with a closed form involving $\zeta(3)$ and $\Li_3$.  The two bases can then be
compared pointwise.  A strictly decreasing universal function
\[
 \Phi(u)=\log_3(3^u-1)-\log_2(2^u-1)
\]
satisfies $\Phi(u)>\log_3 2$ on $(0,1)$.  Consequently every synchronized near-power gap obeys
\[
 |A^d-3^h|>2\,|M^d-2^h|^{\thetaq}
 \qquad(h\in\{\floor{d\beta},\ceil{d\beta}\}),
\]
and the complete gap products have a universal normalized logarithmic surplus
$\kappa N^2+o(N^2)$, where $\kappa=0.7079245703\ldots$.  In particular, after removing the
greatest common divisor of the two complete gap products, the unmatched base-$3$ quotient is
larger than $2^{\binom N2}$ for every $N\ge2$.  This turns the open problem into a precise
second-order valuation-transport question.

A separate algebraic calculation gives the exact norm of $E^k-3^n$ for the canonical radical
$E^d=A/3^a$.  It proves that using all radical conjugates compresses to a single ordinary
binomial gap and cannot, by itself, cross the required scale.  These results are substantial
new structure, but they do not complete the conjecture.
\end{abstract}

\tableofcontents
\clearpage
\pagenumbering{arabic}

\section{Executive summary and claim ledger}
\label{sec:summary}

\subsection{Outcome}

The conjecture remains open in this report.  The work below reaches a new arithmetic object
that was not developed in the supplied unified report: the discriminant of the exact compact
rotation orbit and the integer products obtained by clearing its denominators.  The main gain
is not another equivalent formulation of the original four-logarithm determinant.  It is an
exact decomposition of the ``missing cancellation'' into synchronized integer gaps and a
universal, computable second-order surplus on the base-$3$ side.

\begin{resultbox}
\textbf{Principal new theorem.}  Assume a counterexample and let $M=2^\beta$, $A=3^\beta$ be
its primitive output pair.  Form the complete synchronized gap products
\[
 \mathcal Q_{b,B}(N)=
 \prod_{d=1}^{N-1}
 |B^d-b^{\floor{d\beta}}|^{N-d-C_{d,N}}
 |B^d-b^{\ceil{d\beta}}|^{C_{d,N}},
\]
where $C_{d,N}$ is the exact number of wraps of the length-$(N-d)$ rotation segment by
$\fract{d\beta}$.  Then
\[
 \log_3\mathcal Q_{3,A}(N)
 -\log_2\mathcal Q_{2,M}(N)
 >\binom N2\log_3 2,
\]
and, more precisely, the left side divided by $N^2$ tends to
$\kappa=0.7079245703208558\ldots$.  If
$G_N=\gcd(\mathcal Q_{2,M}(N),\mathcal Q_{3,A}(N))$, then
\[
 \frac{\mathcal Q_{3,A}(N)}{G_N}>2^{\binom N2}.
\]
All three statements are proved in Sections~\ref{sec:amplification} and
\ref{sec:unmatched}.  \statusnew
\end{resultbox}

The final inequality says that a hypothetical counterexample forces at least quadratic
exponential mass in prime powers that occur in the base-$3$ gap product but not in the
base-$2$ gap product.  In the branch $3\nmid A$, this mass is supported entirely on primes
outside the fixed set of divisors of $3A$.  In the branch $3\mid A$, an exact cubic structural
$3$-part can be separated, and the unresolved problem is to control the remaining second-order
transport.

\subsection{Claim ledger}

\begin{center}
\begin{tabularx}{0.98\textwidth}{>{\raggedright\arraybackslash}X p{0.16\textwidth} >{\raggedright\arraybackslash}p{0.31\textwidth}}
\toprule
Statement & Status & Main input \\
\midrule
Primitive generator $\mathrm{Sol}=\Nzero+\Nzero\beta$ and
$M=2^\beta$, $A=3^\beta$ & \statusprior & Supplied unified report \\
Finite-type bound $\norm{q\beta}\ge c q^{-\tau}$ & \statusprior & Baker--W\"ustholz/Matveev as imported in the supplied report \\
Exact wrap and lag factorization of the orbit Vandermonde & \statusnew & Floor arithmetic and a two-case factorization \\
Exact structural-prime denominator and integer numerator product & \statusnew & Unique factorization at $b=2$ or $3$ \\
Orbit-polynomial discriminant identity & \statusnew & Discriminant of a product of integral linear forms \\
Logarithmic-energy limit and closed form $I_b$ & \statusnew & Finite type, discrepancy, uniform integrability, an elementary series integral \\
Pointwise matched-gap amplification & \statusnew & Monotonicity of $b^u/(b^u-1)$ in $b$ \\
Universal surplus $\kappa N^2+o(N^2)$ & \statusnew & Energy limit plus exact integer product identity \\
Quadratic unmatched-prime mass & \statusnew & Normalized surplus and a gcd decomposition \\
Exact canonical-radical norm formula & \statusnew & Cyclotomic multiplicities and $X^D-Y^D$ \\
Complete proof of Alaoglu--Erd\H{o}s & \statusopen & Requires a new cross-base valuation-transport theorem \\
\bottomrule
\end{tabularx}
\end{center}

\section{Scope, imported structure, and non-duplication}
\label{sec:scope}

\subsection{What is imported}

We use two results from the supplied report and do not reprove them.

\begin{enumerate}
\item \statusprior{} If the conjecture is false, there is a unique least nonintegral solution
$\beta>0$ and positive integers
\[
 M=2^\beta,\qquad A=3^\beta
\]
such that every solution is uniquely $n+k\beta$ with $n,k\in\Nzero$.  The pair is affinely
primitive and
\begin{equation}
 \min\{\vpp{2}{M},\vpp{3}{A}\}=0.
 \label{eq:primitive-depth}
\end{equation}
In particular $\beta$ is irrational and every $k\beta$ has nonzero fractional part.

\item \statusprior{} Because $\beta=\log_2 M$ is a ratio of logarithms of fixed algebraic
numbers, a two-logarithm lower bound gives constants $c>0$ and $\tau>0$ with
\begin{equation}
 \norm{q\beta}\ge c q^{-\tau}\qquad(q\ge1),
 \label{eq:finite-type}
\end{equation}
where $\norm{y}$ is distance to the nearest integer.  The constants may be very poor; only
finiteness of $\tau$ is used.
\end{enumerate}

Set throughout
\begin{equation}
 \thetaq=\log_2 3,
 \qquad A=M^{\thetaq}.
 \label{eq:theta}
\end{equation}

\subsection{What is not repeated}

The supplied report already treats the four-exponentials reduction, primitive odd cores,
radical degree, exact counting, block equidistribution, arbitrary-node and semigroup
Vandermonde determinants, $p$-adic assignment minima, exact-grid HCIZ estimates, Loewner and
Smith-profile certificates, polynomial and Pad\'e approximation, compact $S$-integral orbits,
Sturmian and carry languages, natural boundaries, special functions, and twenty-seven
independent continuation reports.  Repackaging any of those arguments would violate the
non-duplication requirement.

The present object is different.  We retain the exact cyclic order of the compact orbit and
multiply \emph{all pairwise differences}.  This creates an integer discriminant after clearing
known denominators, and it allows the base-$2$ and base-$3$ products to be compared one gap at
a time.

\section{Public reconnaissance on the reported AI breakthroughs}
\label{sec:recon}

This section is not used in any proof.  It records what was publicly checkable on August 22,
2026, and separates mathematical verification from attribution.

\subsection{Jacobian conjecture}

A public exact map in dimension three is accompanied by independent symbolic audits and Lean
formalizations.  One convenient form is
\[
\begin{aligned}
 F_1&=(1+xy)^3z+y^2(1+xy)(4+3xy),\\
 F_2&=y+3x(1+xy)^2z+3xy^2(4+3xy),\\
 F_3&=2x-3x^2y-x^3z,
\end{aligned}
\]
with $\det JF=-2$ and an explicit three-point fiber.  The public provenance records credit
Fable for work on the construction; the finite identities can be checked independently.
This refutes the standard conjecture in dimensions at least three.  The plane case remains
open.  See the exact audit~\cite{JacobianAudit}, the Fong--Fable status note~\cite{FongFable},
and a Lean development~\cite{JacobianLean}.  \statusaudit

\subsection{Hadamard matrices}

The public tracker records a coordinated release covering the previously unresolved admissible
orders below $2000$, with exact $\{\pm1\}$ matrices and checker scripts.  The tracker pages did
not present a completely stable review label at the time of access: one page described partial
independent reproduction and review pending, while the later project summary described all
listed orders as verified.  Accordingly, this report treats the result as a strong public
candidate with exact artifacts, not as a substitute for a refereed classification theorem.
See~\cite{HadamardTracker}.  \statusaudit

\subsection{Rank at least thirty}

The public elliptic-curve leaderboard contains an explicit curve, thirty displayed independent
rational points, and an unconditional lower bound $\operatorname{rank}\ge30$ for submission
\#273~\cite{RankThirty}.  The submitter is listed only as ``ranksunbounded.''  I found no
primary public source that identifies that submitter with Fable.  The mathematical artifact is
therefore public; the model attribution in the prompt is not independently verified here.
\statusaudit

\subsection{Alaoglu--Erd\H{o}s}

No indexed public proof, preprint, repository, or Lean development of the claimed new
Alaoglu--Erd\H{o}s route was located in the searches performed for this report.  That is a
search result, not evidence that no private or newly unindexed argument exists.  The work below
was therefore reconstructed independently.

\section{The compact orbit and exact carry calculus}
\label{sec:orbit}

Assume from now on that a counterexample exists.  For $k\ge0$ define
\begin{equation}
 r_k=\floor{k\beta},\qquad t_k=\fract{k\beta},
 \label{eq:rt}
\end{equation}
and, for a base $b>1$,
\begin{equation}
 X_{b,k}=b^{t_k}.
 \label{eq:Xbk}
\end{equation}
For the two arithmetic bases,
\begin{equation}
 X_{2,k}=\frac{M^k}{2^{r_k}},\qquad
 X_{3,k}=\frac{A^k}{3^{r_k}}.
 \label{eq:compact-orbit}
\end{equation}
Thus both coordinates use the same floor sequence and lie in fixed compact intervals.

\subsection{Lag data}

Fix $1\le d<N$.  Put
\begin{equation}
 h_d=\floor{d\beta},\qquad \eta_d=\fract{d\beta},\qquad L_d=N-d.
 \label{eq:lag-data}
\end{equation}
For $0\le i<L_d$, define the carry
\begin{equation}
 \varepsilon_{i,d}=r_{i+d}-r_i-h_d.
 \label{eq:carry}
\end{equation}

\begin{lemma}[Exact two-valued carry]
\label{lem:carry}
For every admissible $i,d$,
\[
 \varepsilon_{i,d}\in\{0,1\},\qquad
 t_{i+d}=t_i+\eta_d-\varepsilon_{i,d}.
\]
Moreover $\varepsilon_{i,d}=1$ if and only if $t_i\ge1-\eta_d$.
\end{lemma}

\begin{proof}
Write $i\beta=r_i+t_i$ and $d\beta=h_d+\eta_d$.  Then
\[
 (i+d)\beta=(r_i+h_d)+(t_i+\eta_d),
\]
where $0<t_i+\eta_d<2$.  Taking the integer and fractional parts gives exactly the stated
formulas.
\end{proof}

Define the exact wrap count
\begin{equation}
 C_{d,N}=\sum_{i=0}^{L_d-1}\varepsilon_{i,d}.
 \label{eq:wrap-count}
\end{equation}
It has the purely integral expression
\begin{equation}
 C_{d,N}=
 \sum_{i=0}^{N-d-1}(r_{i+d}-r_i-h_d).
 \label{eq:wrap-floor-sum}
\end{equation}
No discrepancy estimate is hidden in this identity.

\subsection{Rotation Vandermonde products}

\begin{definition}[Unsquared rotation discriminant]
For $b>1$ and $N\ge2$, set
\begin{equation}
 \Delta_{b,N}=\prod_{0\le i<j<N}|b^{t_j}-b^{t_i}|.
 \label{eq:Delta}
\end{equation}
Its square is the usual Vandermonde discriminant of the $N$ orbit points.
\end{definition}

\begin{theorem}[Exact lag factorization]
\label{thm:lag-factorization}
For every $b>1$, $1\le d<N$, and $L=L_d$,
\begin{align}
 &\prod_{i=0}^{L-1}|b^{t_{i+d}}-b^{t_i}| \\
 &\quad =
 b^{\displaystyle\sum_{i=0}^{L-1}t_i+C_{d,N}(\eta_d-1)}
 (b^{\eta_d}-1)^{L-C_{d,N}}
 (b^{1-\eta_d}-1)^{C_{d,N}}.
 \label{eq:lag-factorization}
\end{align}
Consequently $\Delta_{b,N}$ is the product of the right side over $d=1,\ldots,N-1$.
\end{theorem}

\begin{proof}
If $\varepsilon_{i,d}=0$, Lemma~\ref{lem:carry} gives
\[
 |b^{t_{i+d}}-b^{t_i}|=b^{t_i}(b^{\eta_d}-1).
\]
If $\varepsilon_{i,d}=1$, then $t_{i+d}=t_i+\eta_d-1<t_i$ and
\[
 |b^{t_{i+d}}-b^{t_i}|
 =b^{t_i+\eta_d-1}(b^{1-\eta_d}-1).
\]
Multiplying the $L-C_{d,N}$ nonwrap factors and the $C_{d,N}$ wrap factors gives
\eqref{eq:lag-factorization}.  Every pair $i<j$ occurs uniquely at lag $d=j-i$, proving the
last assertion.
\end{proof}

\begin{remark}
The factors $b^{\eta_d}-1$ and $b^{1-\eta_d}-1$ are a two-branch analogue of a Sudler product.
The exact exponent $C_{d,N}$ retains the carry word instead of replacing it by its expected
value.  This is the feature that later allows an integer gap factorization.
\end{remark}

\section{Exact denominators and integer numerator products}
\label{sec:denominators}

The orbit points are rational when $B=b^\beta$ is an integer.  Their pairwise differences have
completely explicit denominators at the structural prime whenever $b\nmid B$.

\subsection{One pair}

Let $b\in\{2,3\}$, let $B=b^\beta\in\N$, and write
\[
 X_k=\frac{B^k}{b^{r_k}}=b^{t_k}.
\]
For $i<j$,
\begin{equation}
 X_j-X_i=
 \frac{B^j-B^i b^{r_j-r_i}}{b^{r_j}}
 =\frac{B^i(B^{j-i}-b^{r_j-r_i})}{b^{r_j}}.
 \label{eq:pair-rational}
\end{equation}

\begin{theorem}[Exact structural-prime denominator]
\label{thm:exact-denominator}
Assume $b\nmid B$.  Then the numerator in \eqref{eq:pair-rational} is a $b$-adic unit, so the
reduced denominator of $X_j-X_i$ is exactly $b^{r_j}$.
\end{theorem}

\begin{proof}
The imported finite verification implies $\beta>1$; only $\beta>1$ is needed here.  Thus
$r_j-r_i\ge1$.  Modulo $b$, the term $B^{j-i}$ is nonzero while
$b^{r_j-r_i}$ is zero.  Multiplication by the $b$-adic unit $B^i$ preserves nonvanishing.
\end{proof}

\subsection{The complete numerator}

Define
\begin{equation}
 E_N=\sum_{j=1}^{N-1}j r_j
 \label{eq:EN}
\end{equation}
and
\begin{equation}
 P_{b,B}(N)=
 \prod_{0\le i<j<N}|B^j-B^i b^{r_j-r_i}|\in\N.
 \label{eq:P}
\end{equation}

\begin{proposition}[Cleared Vandermonde]
\label{prop:cleared-vdm}
For every $b\in\{2,3\}$ and integer $B=b^\beta$,
\begin{equation}
 P_{b,B}(N)=b^{E_N}\Delta_{b,N}.
 \label{eq:P-Delta}
\end{equation}
If $b\nmid B$, then $P_{b,B}(N)$ is coprime to $b$ and
\[
 \vpp{b}{\Delta_{b,N}}=-E_N.
\]
Moreover
\begin{equation}
 E_N=\frac{\beta}{3}N^3+O_\beta(N^2).
 \label{eq:EN-asymp}
\end{equation}
\end{proposition}

\begin{proof}
Multiplying \eqref{eq:pair-rational} over all pairs contributes $r_j$ once for each
$i=0,\ldots,j-1$, hence the denominator exponent is $E_N$.  The coprimality assertion follows
from Theorem~\ref{thm:exact-denominator}.  Finally,
\[
 E_N=\beta\sum_{j=1}^{N-1}j^2-\sum_{j=1}^{N-1}j t_j,
\]
and the second sum is $O(N^2)$ while $\sum_{j<N}j^2=N^3/3+O(N^2)$.
\end{proof}

\begin{resultbox}
\textbf{Cubic denominator tax.}  The known structural denominator of the complete orbit
Vandermonde has logarithm $\frac13\beta(\log b)N^3+O(N^2)$, whereas the entire archimedean
Vandermonde has logarithm only of order $N^2$ (Theorem~\ref{thm:energy-limit}).  Therefore an
integer-lower-bound argument using only this denominator misses the analytic signal by one full
power of $N$.  This is a different deficit from the one-logarithm node-count deficit in the
supplied report.  \statusnew
\end{resultbox}

\subsection{Gap-only factorization}

The factors $B^i$ in \eqref{eq:P} carry no near-power information.  Remove them by defining
\begin{align}
 \mathcal Q_{b,B}(N)
 &:=\prod_{d=1}^{N-1}
 |B^d-b^{h_d}|^{L_d-C_{d,N}}
 |B^d-b^{h_d+1}|^{C_{d,N}}.
 \label{eq:Q}
\end{align}
Because $d\beta$ is never an integer, every factor is a positive integer.

\begin{theorem}[Exact integer gap product]
\label{thm:gap-product}
For every arithmetic pair $(b,B)$ as above,
\begin{equation}
 P_{b,B}(N)=B^{\binom N3}\mathcal Q_{b,B}(N).
 \label{eq:P-Q}
\end{equation}
\end{theorem}

\begin{proof}
At lag $d$, equation~\eqref{eq:P} contains
\[
 B^i|B^d-b^{h_d+\varepsilon_{i,d}}|
\]
for $i=0,\ldots,L_d-1$.  The wrap count supplies the exponents in
\eqref{eq:Q}.  The total exponent of $B$ is
\[
 \sum_{d=1}^{N-1}\sum_{i=0}^{N-d-1}i
 =\sum_{L=1}^{N-1}\frac{L(L-1)}2
 =\binom N3.
\]
\end{proof}

\section{Orbit polynomials and discriminants}
\label{sec:polynomials}

The same data can be packaged into one integer polynomial.

\begin{definition}[Orbit polynomial]
For $N\ge1$, put
\begin{equation}
 F_{b,B,N}(X)=\prod_{k=0}^{N-1}(b^{r_k}X-B^k)\in\Z[X].
 \label{eq:orbit-polynomial}
\end{equation}
Its roots are the compact orbit points $X_k=B^k/b^{r_k}=b^{t_k}$.
\end{definition}

Let
\begin{equation}
 R_N=\sum_{k=0}^{N-1}r_k,
 \qquad
 S_N=\sum_{i=0}^{N-2}(N-1-i)r_i.
 \label{eq:RS}
\end{equation}

\begin{proposition}[Endpoint data and discriminant]
\label{prop:orbit-disc}
The leading and constant coefficients are
\begin{equation}
 \operatorname{lc}(F_{b,B,N})=b^{R_N},
 \qquad
 |F_{b,B,N}(0)|=B^{\binom N2}.
 \label{eq:endpoints}
\end{equation}
The discriminant is the positive square
\begin{align}
 |\operatorname{Disc}F_{b,B,N}|
 &=\prod_{0\le i<j<N}
 |B^j b^{r_i}-B^i b^{r_j}|^2 \\
 &=b^{2S_N}P_{b,B}(N)^2.
 \label{eq:disc}
\end{align}
\end{proposition}

\begin{proof}
The endpoint formulas are immediate from the product.  For a product of linear forms
$\prod_i(q_iX-p_i)$, direct substitution in the root formula for the discriminant gives
\[
 \operatorname{Disc}\Bigl(\prod_i(q_iX-p_i)\Bigr)
 =\prod_{i<j}(p_iq_j-p_jq_i)^2.
\]
Take $p_i=B^i$ and $q_i=b^{r_i}$.  Each cross determinant contains $b^{r_i}$ times the
corresponding factor in $P_{b,B}(N)$, and $r_i$ occurs $N-1-i$ times.
\end{proof}

For a hypothetical Alaoglu--Erd\H{o}s pair, the two orbit polynomials have identical endpoint
data after normalizing logarithms by the base:
\begin{equation}
 \log_3\operatorname{lc}(F_{3,A,N})
 =\log_2\operatorname{lc}(F_{2,M,N})=R_N,
 \label{eq:lc-match}
\end{equation}
\begin{equation}
 \log_3|F_{3,A,N}(0)|
 =\log_2|F_{2,M,N}(0)|
 =\beta\binom N2.
 \label{eq:ct-match}
\end{equation}
Their discriminants, however, have a universal positive second-order separation.  That is
proved after the logarithmic-energy calculation.

\section{Logarithmic energy of the orbit}
\label{sec:energy}

\subsection{A finite-type uniform-integrability lemma}

For $z_k=t_k=\fract{k\beta}$, the empirical measures
$\mu_N=N^{-1}\sum_{k<N}\delta_{z_k}$ converge weakly to Lebesgue measure on $[0,1]$.
The logarithmic kernel is singular on the diagonal, so weak convergence alone does not justify
passing to pairwise energy.  The finite-type bound \eqref{eq:finite-type} supplies exactly the
missing uniform integrability.

\begin{lemma}[Near-collision tail]
\label{lem:near-tail}
Assume \eqref{eq:finite-type}.  There is $\delta>0$ such that, uniformly for intervals
$I\subset[0,1]$,
\begin{equation}
 \#\{1\le d<N:\fract{d\beta}\in I\}
 =N|I|+O(N^{1-\delta}).
 \label{eq:discrepancy}
\end{equation}
Consequently
\begin{equation}
 \lim_{T\to\infty}\limsup_{N\to\infty}
 \frac1{N^2}\sum_{\substack{0\le i\ne j<N\\|z_i-z_j|<e^{-T}}}
 -\log|z_i-z_j|=0.
 \label{eq:UI}
\end{equation}
\end{lemma}

\begin{proof}
The discrepancy estimate is the standard Erd\H{o}s--Tur\'an consequence of finite type; see
Kuipers--Niederreiter~\cite{KuipersNiederreiter}.  We include the tail reduction because it is
the point at which the hypothesis is used.

Let $A_N(\rho)=\#\{1\le d<N:\norm{d\beta}<\rho\}$.  Equation~\eqref{eq:discrepancy} gives
$A_N(\rho)\le2\rho N+O(N^{1-\delta})$.  If $|z_i-z_j|<\rho$, then
$\norm{(j-i)\beta}<\rho$, and for each lag there are at most $N$ ordered pairs.  The spacing
bound \eqref{eq:finite-type} also gives $|z_i-z_j|\ge cN^{-\tau}$.  Using
$(-\log y-T)_+=\int_T^{-\log y}du$ therefore bounds the normalized tail by
\[
 \int_T^{\tau\log N+O(1)}
 \left(4e^{-u}+O(N^{-\delta})\right)du
 \le4e^{-T}+O(N^{-\delta}\log N).
\]
First let $N\to\infty$, then $T\to\infty$.
\end{proof}

\subsection{Energy limit}

\begin{theorem}[Rotation-discriminant energy]
\label{thm:energy-limit}
For every $b>1$,
\begin{equation}
 \frac{2}{N^2}\log\Delta_{b,N}\longrightarrow I_b,
 \qquad
 I_b=\int_0^1\!\int_0^1\log|b^s-b^t|\,ds\,dt.
 \label{eq:energy-limit}
\end{equation}
The constant has the closed form
\begin{equation}
 I_b=\frac{2\log b}{3}
 -\frac{\pi^2}{3\log b}
 +\frac{2\bigl(\zeta(3)-\Li_3(1/b)\bigr)}{(\log b)^2}.
 \label{eq:Ib}
\end{equation}
\end{theorem}

\begin{proof}
Set
\[
 K_b(s,t)=\log|b^s-b^t|.
\]
The difference
\[
 H_b(s,t)=K_b(s,t)-\log|s-t|
\]
extends continuously to the diagonal, with
$H_b(s,s)=\log((\log b)b^s)$.  Weak convergence of $\mu_N\times\mu_N$ handles $H_b$.  For the
logarithmic term, first truncate at height $-T$, use weak convergence for the bounded continuous
truncation, and then apply Lemma~\ref{lem:near-tail}.  The $N$ omitted diagonal terms contribute
$O(N)/N^2$ after truncation.  This proves \eqref{eq:energy-limit}.

For the integral, let $L=\log b$ and use the density $2(1-u)$ of $u=|s-t|$.  Since
\[
 \log|b^s-b^t|=L\min(s,t)+Lu+\log(1-e^{-Lu}),
\]
the first two expectations total $2L/3$.  Expanding
$\log(1-e^{-Lu})=-\sum_{n\ge1}e^{-nLu}/n$ and integrating termwise gives
\[
 2\int_0^1(1-u)\log(1-e^{-Lu})\,du
 =-\frac{\pi^2}{3L}
 +\frac{2(\zeta(3)-\Li_3(e^{-L}))}{L^2},
\]
which is \eqref{eq:Ib}.
\end{proof}

\subsection{Numerical values}

High-precision evaluation gives
\begin{align}
 I_2&=-1.51660817273459213659\ldots,\label{eq:I2}\\
 I_3&=-0.84829781726239522907\ldots.\label{eq:I3}
\end{align}
These negative values are expected: many orbit points are closer than one in the ordinary
metric, so the product of differences is exponentially small at the $N^2$ scale.

For an implementation check, take $\beta=\log_2 5$.  The following table lists
$2N^{-2}\log\Delta_{b,N}$ and the corresponding normalized base surplus.

\begin{center}
\begin{tabular}{rrrr}
\toprule
$N$ & base $2$ energy & base $3$ energy & normalized surplus $N^{-2}$ \\
\midrule
50   & $-1.4053349594$ & $-0.7474049836$ & $0.6735762212$ \\
100  & $-1.4458192176$ & $-0.7818859890$ & $0.6870864586$ \\
200  & $-1.4834158187$ & $-0.8196667457$ & $0.6970119103$ \\
400  & $-1.4983414907$ & $-0.8318264138$ & $0.7022444033$ \\
800  & $-1.5053190365$ & $-0.8377156849$ & $0.7045973159$ \\
1600 & $-1.5113158350$ & $-0.8433138330$ & $0.7063752646$ \\
3200 & $-1.5133107059$ & $-0.8451095375$ & $0.7069969994$ \\
\midrule
limit & $-1.5166081727$ & $-0.8482978173$ & $0.7079245703$ \\
\bottomrule
\end{tabular}
\end{center}
\statuscomp{}  The table is illustrative only; Theorem~\ref{thm:energy-limit} is proved
independently of the computation.

\section{Matched-gap amplification}
\label{sec:amplification}

This section is the central new arithmetic comparison.

\subsection{The universal one-gap function}

Define for $0<u<1$
\begin{equation}
 \Phi(u)=\log_3(3^u-1)-\log_2(2^u-1).
 \label{eq:Phi}
\end{equation}

\begin{lemma}[Strict gap amplification]
\label{lem:Phi}
The function $\Phi$ is strictly decreasing on $(0,1)$ and
\begin{equation}
 \Phi(u)>\log_3 2\qquad(0<u<1).
 \label{eq:Phi-lower}
\end{equation}
Moreover $\Phi(u)\to+\infty$ as $u\downarrow0$ and
$\Phi(u)\to\log_3 2$ as $u\uparrow1$.
\end{lemma}

\begin{proof}
For $b>1$,
\[
 \frac{d}{du}\log_b(b^u-1)=\frac{b^u}{b^u-1}
 =\frac1{1-b^{-u}}.
\]
For fixed $u>0$ this is strictly decreasing in $b$, hence $\Phi'(u)<0$.  The endpoint at one is
\[
 \log_3(3-1)-\log_2(2-1)=\log_3 2.
\]
The expansion $b^u-1=u\log b+O(u^2)$ proves divergence at zero.
\end{proof}

\subsection{Every synchronized integer gap}

For $d\ge1$, put $h=\floor{d\beta}$ and $\eta=\fract{d\beta}$.  Then
\begin{align}
 A^d-3^h&=3^h(3^\eta-1),\\
 M^d-2^h&=2^h(2^\eta-1).
\end{align}
For the ceiling branch, the absolute gaps are
$3^h(3-3^\eta)$ and $2^h(2-2^\eta)$.

\begin{theorem}[Pointwise matched-gap inequality]
\label{thm:pointwise-gap}
For every $d\ge1$ and
$h\in\{\floor{d\beta},\ceil{d\beta}\}$,
\begin{equation}
 |A^d-3^h|>2\,|M^d-2^h|^{\thetaq}.
 \label{eq:pointwise-gap}
\end{equation}
Equivalently,
\begin{equation}
 \log_3|A^d-3^h|-\log_2|M^d-2^h|>\log_3 2.
 \label{eq:pointwise-log}
\end{equation}
\end{theorem}

\begin{proof}
For $h=\floor{d\beta}$, the $h$ terms cancel in the normalized logarithms and the left side of
\eqref{eq:pointwise-log} is $\Phi(\eta)$.  For $h=\ceil{d\beta}$ it is
\[
 \log_3(3-3^\eta)-\log_2(2-2^\eta)=\Phi(1-\eta),
\]
because $b-b^\eta=b^\eta(b^{1-\eta}-1)$.  Lemma~\ref{lem:Phi} proves both cases.
Multiplying \eqref{eq:pointwise-log} by $\log3$ and using
$\thetaq=\log_2 3$ gives \eqref{eq:pointwise-gap}.
\end{proof}

\subsection{Complete products}

\begin{theorem}[Finite universal surplus]
\label{thm:finite-surplus}
For every $N\ge2$,
\begin{equation}
 \mathcal G_N:=
 \log_3\mathcal Q_{3,A}(N)
 -\log_2\mathcal Q_{2,M}(N)
 >\binom N2\log_3 2.
 \label{eq:G-lower}
\end{equation}
Equivalently,
\begin{equation}
 \mathcal Q_{3,A}(N)>
 2^{\binom N2}\,\mathcal Q_{2,M}(N)^{\thetaq}.
 \label{eq:Q-product-ineq}
\end{equation}
The same surplus can be written in any of the following exact forms:
\begin{align}
 \mathcal G_N
 &=\log_3 P_{3,A}(N)
   -\log_2 P_{2,M}(N),
 \label{eq:G-P}\\
 &=\log_3\Delta_{3,N}
   -\log_2\Delta_{2,N},
 \label{eq:G-Delta}\\
 &=\sum_{0\le i<j<N}\Phi(|t_j-t_i|),
 \label{eq:G-pairwise}\\
 &=\sum_{d=1}^{N-1}
 \bigl((L_d-C_{d,N})\Phi(\eta_d)
       +C_{d,N}\Phi(1-\eta_d)\bigr).
 \label{eq:G-lag}
\end{align}
\end{theorem}

\begin{proof}
Apply Theorem~\ref{thm:pointwise-gap} to each of the
$(L_d-C_{d,N})$ floor gaps and $C_{d,N}$ ceiling gaps at lag $d$.  The total number of factors
is $\sum_{d=1}^{N-1}L_d=\binom N2$, giving \eqref{eq:G-lower} and
\eqref{eq:Q-product-ineq}.

Equation~\eqref{eq:G-P} follows from Theorem~\ref{thm:gap-product}: the monomial terms cancel
because
\[
 \log_3 A=\log_2 M=\beta.
\]
Equation~\eqref{eq:G-Delta} follows from Proposition~\ref{prop:cleared-vdm}, since the common
$E_N$ cancels.  For a pair with $s=\min(t_i,t_j)$ and $u=|t_i-t_j|$,
\[
 \log_b|b^{t_j}-b^{t_i}|=s+\log_b(b^u-1),
\]
so the $s$ terms cancel between bases and leave $\Phi(u)$.  This proves
\eqref{eq:G-pairwise}; grouping by lags gives \eqref{eq:G-lag}.
\end{proof}

\subsection{The exact second-order constant}

\begin{theorem}[Universal matched-gap energy]
\label{thm:kappa}
One has
\begin{equation}
 \frac{\mathcal G_N}{N^2}\longrightarrow\kappa,
 \qquad
 \kappa=\frac12\left(\frac{I_3}{\log3}-\frac{I_2}{\log2}\right),
 \label{eq:kappa}
\end{equation}
where
\begin{align}
 \kappa
 &=\frac{\pi^2}{6}
 \left(\frac1{(\log2)^2}-\frac1{(\log3)^2}\right) \\
 &\quad +\frac{\zeta(3)-\Li_3(1/3)}{(\log3)^3}
 -\frac{\zeta(3)-\Li_3(1/2)}{(\log2)^3} \\
 &=0.7079245703208558078744\ldots.
 \label{eq:kappa-closed}
\end{align}
In particular
\begin{equation}
 \kappa>\frac12\log_3 2=0.3154648767\ldots.
 \label{eq:kappa-lower}
\end{equation}
\end{theorem}

\begin{proof}
Combine \eqref{eq:G-Delta} with Theorem~\ref{thm:energy-limit}.  Substitution of
\eqref{eq:Ib} gives the closed form; the $2/3$ terms cancel.  Inequality
\eqref{eq:kappa-lower} also follows without numerical evaluation by dividing
\eqref{eq:G-lower} by $N^2$ and taking the limit.
\end{proof}

\begin{corollary}[Discriminant separation]
\label{cor:disc-separation}
The orbit polynomials satisfy
\begin{equation}
 \log_3|\operatorname{Disc}F_{3,A,N}|
 -\log_2|\operatorname{Disc}F_{2,M,N}|
 =2\mathcal G_N
 >2\binom N2\log_3 2,
 \label{eq:disc-separation}
\end{equation}
while their normalized leading and constant coefficients agree exactly by
\eqref{eq:lc-match}--\eqref{eq:ct-match}.  The left side divided by $N^2$ tends to $2\kappa$.
\end{corollary}

\begin{proof}
Use Proposition~\ref{prop:orbit-disc}.  The common $S_N$ cancels, leaving twice
\eqref{eq:G-P}.
\end{proof}

\section{Quadratic unmatched-prime mass}
\label{sec:unmatched}

The universal surplus has an immediate arithmetic interpretation that does not mention real
powers.

\begin{theorem}[Unmatched base-$3$ quotient]
\label{thm:unmatched}
Let
\begin{equation}
 H_N=\gcd(\mathcal Q_{2,M}(N),\mathcal Q_{3,A}(N)),
 \qquad
 Y_N=\frac{\mathcal Q_{3,A}(N)}{H_N}.
 \label{eq:HN-YN}
\end{equation}
Then for every $N\ge2$,
\begin{equation}
 Y_N>2^{\binom N2}.
 \label{eq:unmatched-finite}
\end{equation}
More precisely,
\begin{equation}
 \liminf_{N\to\infty}\frac{\log Y_N}{N^2}
 \ge\kappa\log3
 =0.7777346324045835\ldots.
 \label{eq:unmatched-asymp}
\end{equation}
\end{theorem}

\begin{proof}
Write $\mathcal Q_{2,M}=H_NX_N$ and $\mathcal Q_{3,A}=H_NY_N$ with
$\gcd(X_N,Y_N)=1$.  Since $\log3>\log2$,
\begin{align*}
 \mathcal G_N
 &=\log H_N\left(\frac1{\log3}-\frac1{\log2}\right)
 +\log_3 Y_N-\log_2 X_N\\
 &\le\log_3 Y_N.
\end{align*}
Apply Theorems~\ref{thm:finite-surplus} and~\ref{thm:kappa}.
\end{proof}

This is stronger than saying that the two gap products are unequal.  It forces a rapidly
growing quotient supported on prime powers absent from the dyadic gap product.

\subsection{Support in the triadically primitive branch}

Write
\begin{equation}
 A=3^bV,\qquad b=\vpp{3}{A},\qquad 3\nmid V.
 \label{eq:Adepth}
\end{equation}

\begin{corollary}[External support when $3\nmid A$]
\label{cor:external-support}
If $b=0$, then
\begin{equation}
 \gcd(\mathcal Q_{3,A}(N),3A)=1.
 \label{eq:Q3-coprime}
\end{equation}
Consequently the quotient $Y_N$ in Theorem~\ref{thm:unmatched} is supported on primes outside
the fixed set of divisors of $3A$ and has logarithmic size at least
$(\kappa\log3+o(1))N^2$.
\end{corollary}

\begin{proof}
Every factor of $\mathcal Q_{3,A}$ is $|A^d-3^h|$ with $h\ge1$.  Modulo $3$ it is nonzero
because $3\nmid A$.  If $p\mid A$ and $p\ne3$, then $A^d-3^h\equiv-3^h\not\equiv0\pmod p$.
\end{proof}

\subsection{Exact structural part when \texorpdfstring{$3\mid A$}{3 divides A}}

If $b>0$, the $3$-part of the triadic gap product is cubic and must be separated before the
quadratic surplus can be interpreted as genuinely new primes.  Put
\[
 \gamma=\beta-b=\log_3V>0.
\]
Because $\gamma$ is irrational, $\ceil{d\gamma}\ge1$ for every $d\ge1$.

\begin{proposition}[Exact structural $3$-part]
\label{prop:structural-three}
If $b>0$, then
\begin{equation}
 \vpp{3}{\mathcal Q_{3,A}(N)}
 =b\binom{N+1}{3}
 +\sum_{\substack{1\le d<N\\d\gamma<1}}
 (N-d-C_{d,N})\,\vpp{3}{V^d-1}.
 \label{eq:structural-three}
\end{equation}
The sum on the right has finitely many $d$, independent of $N$, and is $O_A(N)$.  Hence
\begin{equation}
 \vpp{3}{\mathcal Q_{3,A}(N)}
 =\frac{b}{6}N^3+O_A(N^2).
 \label{eq:structural-three-asymp}
\end{equation}
\end{proposition}

\begin{proof}
For $h=\ceil{d\beta}=bd+\ceil{d\gamma}>bd$,
\[
 \vpp{3}{A^d-3^h}=bd.
\]
For $h=\floor{d\beta}=bd+\floor{d\gamma}$ the same holds unless
$\floor{d\gamma}=0$.  In that exceptional case the valuation is
$bd+\vpp{3}{V^d-1}$.  The floor branch appears with exponent
$N-d-C_{d,N}$.  Finally
\[
 \sum_{d=1}^{N-1}d(N-d)=\binom{N+1}{3}.
\]
\end{proof}

There is a symmetric statement for $a=\vpp{2}{M}>0$.  The primitive-depth condition
\eqref{eq:primitive-depth} says that the two structural cubic terms never occur simultaneously.
The open problem is not to discover these terms; they are exact.  It is to transport the
remaining prime mass across the two bases at the $N^2$ scale.

\subsection{The valuation-transport target}

Expanding the universal surplus into prime valuations gives the exact identity
\begin{equation}
 \mathcal G_N=
 \sum_p\left(
 \frac{\vpp{p}{\mathcal Q_{3,A}(N)}}{\log3}
 -\frac{\vpp{p}{\mathcal Q_{2,M}(N)}}{\log2}
 \right)\log p.
 \label{eq:valuation-identity}
\end{equation}
This is merely unique factorization, but it pinpoints what a closing theorem must do.

\begin{criterion}[Cross-base valuation transport]
\label{crit:transport}
Any arithmetic theorem, applicable to every primitive integer pair $(M,A)$, which forces
\begin{equation}
 \limsup_{N\to\infty}\frac1{N^2}
 \sum_p\left(
 \frac{\vpp{p}{\mathcal Q_{3,A}(N)}}{\log3}
 -\frac{\vpp{p}{\mathcal Q_{2,M}(N)}}{\log2}
 \right)\log p<\kappa
 \label{eq:transport-target}
\end{equation}
would prove the Alaoglu--Erd\H{o}s conjecture.
\end{criterion}

\begin{proof}
Under a counterexample, Theorem~\ref{thm:kappa} says that the same expression tends to
$\kappa$.
\end{proof}

Criterion~\ref{crit:transport} is intentionally phrased as a sufficient condition, not as a
new name for the conjecture.  Its value is that every summand concerns valuations of explicit
ordinary integer differences
\[
 M^d-2^{\floor{d\beta}},\quad M^d-2^{\ceil{d\beta}},\quad
 A^d-3^{\floor{d\beta}},\quad A^d-3^{\ceil{d\beta}}.
\]
There are no logarithmic coefficients inside those valuations.  The transcendental relation
enters only through the common Beatty exponents.

\begin{warning}
A small-gcd theorem by itself does not close the problem.  Standard subspace-theorem estimates
for common divisors tend to make $H_N$ smaller, which makes the unmatched quotient in
Theorem~\ref{thm:unmatched} larger.  The needed input is a cross-base \emph{upper} control on
the normalized unshared mass, or a mechanism that converts its growth into a forbidden third
algebraic exponential.
\end{warning}

\section{Canonical-radical norm compression}
\label{sec:radical}

The supplied report reduces a counterexample to a canonical radical.  We now compute exactly
what all of its conjugates contribute to a near-power gap.

\subsection{Setup}

\statusprior{}  There are integers $w>1$ odd, $d\ge1$, $a\ge0$, and $A\ge1$ such that
\begin{equation}
 E=w^{\thetaq}>0,\qquad E^d=q=\frac{A}{3^a},
 \label{eq:E}
\end{equation}
$d$ is the least positive exponent for which $E^d\in\Q$, and
$X^d-q$ is irreducible over $\Q$.

Fix $k\ge1$ and $n\ge0$.  Put
\begin{equation}
 g=\gcd(k,d),\qquad K=\frac{k}{g},\qquad D=\frac{d}{g}.
 \label{eq:gKD}
\end{equation}

\begin{theorem}[Exact cyclotomic norm]
\label{thm:radical-norm}
The field norm is
\begin{align}
 \left|N_{\Q(E)/\Q}(E^k-3^n)\right|
 &=\left|q^K-3^{nD}\right|^g \\
 &=\frac{|A^K-3^{aK+nD}|^g}{3^{ak}}.
 \label{eq:radical-norm}
\end{align}
The numerator is nonzero.  Furthermore
\begin{equation}
 |E^k-3^n|
 =\frac{|A^K-3^{aK+nD}|}
 {3^{aK}\displaystyle\sum_{j=0}^{D-1}(E^k)^{D-1-j}(3^n)^j},
 \label{eq:radical-factor}
\end{equation}
and hence
\begin{equation}
 |E^k-3^n|
 \ge\frac{3^{-aK}}{D\max(E^k,3^n)^{D-1}}.
 \label{eq:radical-lower-elementary}
\end{equation}
\end{theorem}

\begin{proof}
The conjugates of $E$ are $\zeta_d^jE$.  The multiset of $k$th powers
$\zeta_d^{jk}$ contains every $D$th root of unity exactly $g$ times.  Therefore
\[
 \prod_{j=0}^{d-1}\bigl((\zeta_d^jE)^k-3^n\bigr)
 =\pm\bigl((E^k)^D-3^{nD}\bigr)^g
 =\pm(q^K-3^{nD})^g.
\]
Substitution of $q=A/3^a$ gives \eqref{eq:radical-norm}.  If its numerator vanished, positivity
would give $E^k=3^n$, and then
$k\thetaq\log w=n\log3$, so $w^k=2^n$, impossible for odd $w>1$.

Finally use
\[
 X^D-Y^D=(X-Y)\sum_{j=0}^{D-1}X^{D-1-j}Y^j
\]
with $X=E^k$, $Y=3^n$, and clear $3^{aK}$.  The denominator sum is at most
$D\max(X,Y)^{D-1}$.
\end{proof}

\subsection{Baker-strengthened relative separation}

Because $E$ and $3$ are multiplicatively independent positive algebraic numbers, Baker's
theorem gives constants $C,\mu>0$ such that
\begin{equation}
 |k\log E-n\log3|\ge C(k+n)^{-\mu}
 \label{eq:Baker-E}
\end{equation}
whenever the form is nonzero.  Standard comparison with $e^z-1$ yields
\begin{equation}
 |E^k-3^n|
 \ge C'\max(E^k,3^n)(k+n)^{-\mu'}
 \label{eq:relative-radical}
\end{equation}
for suitable $C',\mu'>0$, after harmless adjustment outside the close regime.

This is stronger than \eqref{eq:radical-lower-elementary}, but it still matches the best
rational-return scale rather than contradicting it.  The norm has not created a new algebraic
quantity; it has compressed all conjugates to the ordinary binomial difference
$A^K-3^{aK+nD}$.

\begin{resultbox}
\textbf{Exact no-go statement.}  Any argument that uses the conjugates of $E$ only through the
field norm of $E^k-3^n$ reduces exactly to one two-logarithm gap.  Norm amplification alone
cannot create the missing third independent exponential or a super-polynomial separation.
A successful radical argument must retain more than the total product of conjugates---for
example, a signed subset product, a resolvent sensitive to the real embedding, or a coupled
finite-place invariant.  \statusnew
\end{resultbox}

\section{Why the new identities still stop short}
\label{sec:failure}

\subsection{The integer lower bound has the wrong scale}

When $b\nmid B$, the rational number $\Delta_{b,N}$ has exact denominator
$b^{E_N}$, so nonvanishing gives only
\[
 |\Delta_{b,N}|\ge b^{-E_N}
 =\exp\left(-\frac{\beta\log b}{3}N^3+O(N^2)\right).
\]
The true size is
\[
 |\Delta_{b,N}|=
 \exp\left(\frac{I_b}{2}N^2+o(N^2)\right).
\]
The lower bound is therefore exponentially weaker by one power of $N$.  Taking squares,
discriminants, or products of the same denominators cannot change that order.

\subsection{The universal surplus is compatible with ordinary arithmetic}

The inequality
\[
 \mathcal Q_{3,A}>2^{\binom N2}\mathcal Q_{2,M}^{\thetaq}
\]
is strong but points in a direction that large-base geometry naturally permits.  Both gap
products themselves have logarithm of order $N^3$; the surplus is only their second-order
difference.  An arbitrary pair of unrelated integer products can easily have such a surplus.
The conjecture can be reached only by exploiting the fact that both products use the
\emph{same} Beatty exponents and arise from the same primitive output pair.

\subsection{Why a gcd upper bound is not enough}

Theorem~\ref{thm:unmatched} already says that the quotient remaining after the gcd is enormous.
A theorem of Bugeaud--Corvaja--Zannier type~\cite{BCZ} that makes the common divisor even smaller
does not contradict anything.  What is needed is one of the following stronger mechanisms:

\begin{enumerate}
\item an upper bound for the weighted unmatched base-$3$ prime mass at the exact $N^2$ scale;
\item a rule transporting a fixed positive proportion of that mass into the dyadic gap product;
\item a primitive-divisor theorem whose new primes manufacture a third algebraic exponential;
\item a cross-resultant or resolvent whose product formula turns the surplus into an upper, not
      a lower, archimedean bound.
\end{enumerate}

\subsection{Why the radical norm is not the missing transport}

Equation~\eqref{eq:radical-norm} has only one numerator, and Baker controls it polynomially in
the exponent.  It does not connect primes dividing $M^d-2^h$ to primes dividing $A^d-3^h$.
Thus the radical and discriminant routes currently meet only at the same unresolved transport
problem.

\section{Concrete research program}
\label{sec:program}

\subsection{Priority A: cross-base local orders}

For a prime $p\nmid6MA$, simultaneous divisibility
\[
 p\mid M^d-2^h,
 \qquad
 p\mid A^d-3^h
\]
means that the reduction of $(M,A)^d(2,3)^{-h}$ is the identity in
$(\mathbb F_p^\times)^2$.  Fixing a generator of $\mathbb F_p^\times$ turns this into a
$2\times2$ modular linear system.  The determinant of its discrete-log matrix is the finite
place analogue of the real logarithmic determinant.  A useful theorem would compare the
orders of this kernel for many $p$ with the exact Beatty set
$h\in\{\floor{d\beta},\ceil{d\beta}\}$.

The key requirement is quantitative: an $o(N^2)$ estimate is not enough if it concerns only the
number of primes.  It must control the log-weighted valuations in
\eqref{eq:valuation-identity}.

\subsection{Priority B: primitive divisors of moving-exponent gaps}

Classical Zsigmondy theory handles $a^n-b^n$.  Here the relevant factors are
\[
 A^d-3^{\floor{d\beta}},\qquad
 A^d-3^{\ceil{d\beta}},
\]
with two exponents moving along an irrational line.  A theorem guaranteeing a primitive prime
for most $d$, together with a criterion excluding that prime from the corresponding dyadic
gap, would turn Theorem~\ref{thm:unmatched} from aggregate mass into a supply of genuinely new
prime directions.  The second step is essential; primitive divisors on the triadic side alone
do not force another algebraic exponential.

\subsection{Priority C: signed or partial radical norms}

The full norm destroys the real embedding.  A more promising radical object is a partial
product over a Galois-stable but not full set of conjugates, or a resolvent that records which
conjugate lies on the positive real branch.  Its denominator and local valuations must still be
controlled exactly.  Equation~\eqref{eq:radical-norm} is a test: any proposed invariant that
collapses to that formula has not added information.

\subsection{Priority D: orbit-polynomial comparison}

The polynomials $F_{2,M,N}$ and $F_{3,A,N}$ have equal normalized leading and constant
coefficients but discriminants separated by $2\kappa N^2+o(N^2)$.  A cross-base coefficient
inequality sensitive to complete splitting over $\Q$, rather than to Mahler measure alone,
could convert this into a contradiction.  Standard discriminant--Mahler bounds are too coarse:
their leading scale is $N^3$ and they tolerate the observed $N^2$ difference.

\subsection{Priority E: formalization before further optimization}

The finite identities in Sections~\ref{sec:orbit}, \ref{sec:denominators},
\ref{sec:polynomials}, \ref{sec:amplification}, and \ref{sec:radical} are suitable for immediate
Lean formalization.  Doing so would protect the exponent bookkeeping---especially the wrap
multiplicities and discriminant powers---before a deeper valuation experiment is built on top
of them.

\section{Lean-oriented formalization blueprint}
\label{sec:lean}

No claim of kernel verification is made for this continuation.  The following decomposition
keeps all analytic inputs explicit.

\subsection{Module \texttt{RotationCarry}}

Definitions:
\[
 r(k)=\floor{k\beta},\quad t(k)=k\beta-r(k),\quad
 \varepsilon(i,d)=r(i+d)-r(i)-r(d).
\]
Targets:
\begin{itemize}
\item $\varepsilon(i,d)=0\lor\varepsilon(i,d)=1$;
\item the exact fractional-part update;
\item the floor-sum formula for $C_{d,N}$.
\end{itemize}
All are elementary ordered-ring/floor lemmas.

\subsection{Module \texttt{RotationVandermonde}}

Define the finite product $\Delta_{b,N}$ and prove Theorem~\ref{thm:lag-factorization} by
partitioning the finite index set according to the carry.  No transcendence theorem is needed;
irrationality of $\beta$ is used only to make every factor nonzero.

\subsection{Module \texttt{OrbitGapProduct}}

For prime $b$ and natural $B$, define the integer numerator matrix
\[
 N_{ij}=B^j-B^i b^{r_j-r_i}.
\]
Formalize:
\begin{itemize}
\item the exact denominator under $b\nmid B$;
\item $P=b^{E_N}\Delta$;
\item $P=B^{\binom N3}\mathcal Q$;
\item the finite sums $\sum_{j<N}jr_j$ and
      $\sum_{d<N}\sum_{i<N-d}i=\binom N3$.
\end{itemize}

\subsection{Module \texttt{OrbitPolynomialDiscriminant}}

Formalize the general identity
\[
 \operatorname{Disc}\left(\prod_i(q_iX-p_i)\right)
 =\prod_{i<j}(p_iq_j-p_jq_i)^2
\]
for a field, then specialize to integral coefficients.  This module should be independent of
real powers once the integer sequences are supplied.

\subsection{Module \texttt{MatchedGapAmplification}}

The key analytic lemma is one-dimensional:
\[
 u\mapsto\log_3(3^u-1)-\log_2(2^u-1)
\]
is strictly decreasing.  Mathlib's derivative rules for real powers and logarithms should make
the proof direct.  The integer product inequality and unmatched-gcd corollary are then finite
algebra and order arithmetic.

A schematic endpoint theorem is:
\begin{verbatim}
theorem unmatched_triadic_gap_mass
    (hpair : A = Real.rpow M (Real.logb 2 3))
    (hbeta : ... exact floor synchronization ... ) :
    2 ^ Nat.choose N 2 < Q3 N / Nat.gcd (Q2 N) (Q3 N)
\end{verbatim}
The actual statement should avoid natural-number division by expressing the quotient through a
witness from the gcd divisibility theorem.

\subsection{Module \texttt{RotationLogEnergy}}

Separate the theorem into two interfaces:
\begin{enumerate}
\item a finite-type discrepancy hypothesis implying uniform integrability of the logarithmic
kernel;
\item the elementary evaluation of $I_b$.
\end{enumerate}
The finite algebraic conclusions do not depend on this module; only the limit constant
$\kappa$ does.

\subsection{Module \texttt{CanonicalRadicalNorm}}

Formalize the multiset map $j\mapsto jk\pmod d$, its fibers of size $g=\gcd(k,d)$, and the
cyclotomic product.  The exact norm theorem can be stated first for an abstract root of
$X^d-q$ and then specialized to the positive real radical.

\section{Computational reconnaissance}
\label{sec:computation}

The calculations in this report are symbolic.  Computation was used only to check indexing and
to evaluate constants.  The following self-contained Python script verifies the finite product
factorization for exact integers and reproduces the energy table.

\begin{lstlisting}[style=pythonstyle,caption={Exact product checks and numerical energy.}]
from __future__ import annotations

from math import comb, floor, log
from typing import List

import mpmath as mp
import numpy as np

mp.mp.dps = 100


def floor_sequence(base: int, output: int, n: int) -> List[int]:
    """r_k=floor(k log_base(output)), evaluated at high precision."""
    beta = mp.log(output) / mp.log(base)
    return [int(mp.floor(k * beta)) for k in range(n)]


def direct_numerator_product(base: int, output: int, n: int) -> int:
    r = floor_sequence(base, output, n)
    ans = 1
    for i in range(n):
        for j in range(i + 1, n):
            ans *= abs(output**j - output**i * base ** (r[j] - r[i]))
    return ans


def grouped_gap_product(base: int, output: int, n: int) -> int:
    r = floor_sequence(base, output, n)
    q = 1
    for d in range(1, n):
        h = r[d]
        wraps = sum(r[i + d] - r[i] - h for i in range(n - d))
        q *= abs(output**d - base**h) ** (n - d - wraps)
        q *= abs(output**d - base ** (h + 1)) ** wraps
    return output ** comb(n, 3) * q


def verify_exact(base: int, output: int, n: int) -> None:
    assert output % base != 0
    left = direct_numerator_product(base, output, n)
    right = grouped_gap_product(base, output, n)
    assert left == right

    r = floor_sequence(base, output, n)
    for i in range(n):
        for j in range(i + 1, n):
            numerator = output**j - output**i * base ** (r[j] - r[i])
            assert numerator % base != 0  # exact reduced denominator


def I(base: int) -> mp.mpf:
    L = mp.log(base)
    return (2 * L / 3 - mp.pi**2 / (3 * L)
            + 2 * (mp.zeta(3) - mp.polylog(3, mp.mpf(1) / base)) / L**2)


def empirical_energy(beta: float, base: int, n: int) -> float:
    t = (np.arange(n, dtype=np.float64) * beta) % 1.0
    x = np.power(float(base), t)
    total = 0.0
    for i in range(n):
        total += np.log(np.abs(x[i + 1:] - x[i])).sum()
    return 2.0 * total / (n * n)


if __name__ == "__main__":
    verify_exact(2, 5, 14)
    verify_exact(3, 10, 14)

    beta = log(5, 2)
    for n in [50, 100, 200, 400, 800, 1600, 3200]:
        e2 = empirical_energy(beta, 2, n)
        e3 = empirical_energy(beta, 3, n)
        surplus = 0.5 * (e3 / log(3) - e2 / log(2))
        print(n, e2, e3, surplus)

    kappa = mp.mpf("0.5") * (I(3) / mp.log(3) - I(2) / mp.log(2))
    print("I2 =", I(2))
    print("I3 =", I(3))
    print("kappa =", kappa)
\end{lstlisting}

The exact check uses $(b,B)=(2,5)$ and $(3,10)$ separately; no counterexample is assumed or
simulated.  The energy test can use any finite-type irrational rotation because the limiting
constant depends on $b$, not on the rotation angle.

\section{Conclusions}
\label{sec:conclusion}

The attached report identified two recurring obstacles: the $2\times2$ four-exponentials wall
and a shortage of one logarithmic factor in determinant arguments.  The present continuation
does not remove either wall, but it changes the arithmetic interface in three concrete ways.

First, every compact-orbit Vandermonde is now factored exactly by lag and carry, and its
structural denominator is known with equality.  This proves that denominator clearing by itself
has a cubic tax against a quadratic analytic signal.

Second, comparing the two bases before discarding the carry structure produces a universal
pointwise inequality.  Its complete product is not merely nonzero: it has a fixed quadratic
surplus and forces an exponentially large quotient of triadic gap prime powers not shared by
the dyadic gaps.  The constant $\kappa$ is explicit and independent of the hypothetical
counterexample.

Third, the full conjugate norm of the canonical radical is computed exactly and shown to
collapse to one ordinary binomial gap.  This eliminates a broad family of norm-only arguments
and specifies what an improved radical invariant must retain.

The clearest next theorem is therefore not another reformulation of
$\log M\log3=\log A\log2$.  It is a second-order valuation-transport estimate for the explicit
moving-exponent gaps in \eqref{eq:valuation-identity}.  Proving such an estimate with a constant
strictly below $\kappa$, or converting the forced unmatched primes into a third algebraic
exponential, would complete the conjecture.

\appendix

\section{Details of the energy integral}
\label{app:integral}

For completeness, let $L=\log b$.  Symmetry and the change of variables $u=|s-t|$ give
\begin{align*}
 I_b
 &=2\int_0^1\int_0^{1-u}
 \left(Ls+\log(e^{Lu}-1)\right)ds\,du\\
 &=\frac{L}{3}+2\int_0^1(1-u)\log(e^{Lu}-1)\,du\\
 &=\frac{2L}{3}+2\int_0^1(1-u)\log(1-e^{-Lu})\,du.
\end{align*}
Now
\[
 \int_0^1(1-u)e^{-nLu}\,du
 =\frac1{nL}-\frac{1-e^{-nL}}{n^2L^2}.
\]
Therefore
\begin{align*}
 2\int_0^1(1-u)\log(1-e^{-Lu})\,du
 &=-2\sum_{n\ge1}\frac1n
 \left(\frac1{nL}-\frac{1-e^{-nL}}{n^2L^2}\right)\\
 &=-\frac{2\zeta(2)}L
 +\frac{2(\zeta(3)-\Li_3(e^{-L}))}{L^2},
\end{align*}
which is \eqref{eq:Ib}.

\section{A compact dependency graph}
\label{app:dependencies}

\begin{center}
\begin{tabularx}{0.98\textwidth}{>{\raggedright\arraybackslash}p{0.28\textwidth} >{\raggedright\arraybackslash}X}
\toprule
Result & Dependencies \\
\midrule
Exact carry and lag factorization & floor addition; no imported transcendence \\
Exact denominator & $b$ prime, $b\nmid B$, $\beta>1$ \\
Integer gap product & exact carry and finite product reindexing \\
Orbit discriminant & general discriminant of integral linear factors \\
Energy limit & irrational rotation, finite-type bound, discrepancy \\
Closed form $I_b$ & elementary integration and power series \\
Pointwise amplification & $A=M^{\thetaq}$ and monotonicity of $\Phi$ \\
Finite surplus & pointwise amplification and exact wrap multiplicities \\
Limit $\kappa$ & finite surplus identity plus energy theorem \\
Unmatched quotient & finite surplus plus gcd factorization \\
Radical norm & irreducible binomial and cyclotomic multiplicity \\
\bottomrule
\end{tabularx}
\end{center}

\begin{thebibliography}{99}
\small

\bibitem{AE1944}
L.~Alaoglu and P.~Erd\H{o}s,
\newblock \emph{On highly composite and similar numbers},
\newblock Trans. Amer. Math. Soc. \textbf{56} (1944), 448--469.

\bibitem{Baker}
A.~Baker,
\newblock \emph{Transcendental Number Theory},
\newblock Cambridge University Press, 1975.

\bibitem{BakerWustholz}
A.~Baker and G.~W\"ustholz,
\newblock \emph{Logarithmic Forms and Diophantine Geometry},
\newblock Cambridge University Press, 2007.

\bibitem{Matveev}
E.~M.~Matveev,
\newblock An explicit lower bound for a homogeneous rational linear form in logarithms of algebraic numbers II,
\newblock Izv. Math. \textbf{64} (2000), 1217--1269.

\bibitem{KuipersNiederreiter}
L.~Kuipers and H.~Niederreiter,
\newblock \emph{Uniform Distribution of Sequences},
\newblock Wiley, 1974.

\bibitem{BCZ}
Y.~Bugeaud, P.~Corvaja, and U.~Zannier,
\newblock An upper bound for the G.C.D. of $a^n-1$ and $b^n-1$,
\newblock Math. Z. \textbf{243} (2003), 79--84.

\bibitem{UnifiedReport}
The ProveIt project,
\newblock \emph{The Alaoglu--Erd\H{o}s Integer-Exponent Conjecture: Unified Report},
\newblock supplied source document, August 21, 2026.

\bibitem{JacobianAudit}
B.~Naskr\k{e}cki,
\newblock \emph{A three-dimensional Keller map with a three-point fiber: exact audit and structural analysis},
\newblock public repository, 2026,
\newblock \url{https://github.com/nasqret/jacobian-counterexample}.

\bibitem{FongFable}
D.~Fong and Fable,
\newblock \emph{The State of the Jacobian Conjectures},
\newblock status note, July 23, 2026,
\newblock \url{https://jacobianconjectures.com/jacobian/note/}.

\bibitem{JacobianLean}
D.~Cureton,
\newblock \emph{Levent Alp\"oge/Fable 5's counterexample to the Jacobian conjecture},
\newblock Lean 4 development, 2026,
\newblock \url{https://github.com/deancureton/jacobian/blob/main/Jacobian/Counterexample.lean}.

\bibitem{HadamardTracker}
VibeMathed contributors,
\newblock \emph{Hadamard Matrix Conjecture---Order 668} and project status records,
\newblock accessed August 22, 2026,
\newblock \url{https://vibemathed.com/problem/hadamard-matrix-order-668}.

\bibitem{RankThirty}
Elliptic Curve Rank Leaderboard,
\newblock \emph{Curve \#273}, submitted August 20, 2026,
\newblock \url{https://elliptic-rank.icarm.cloud/curve/273}.

\end{thebibliography}

\end{document}
